\documentclass[../../thesis_main.tex]{subfiles}
\ifSubfilesClassLoaded{\externaldocument{../../thesis_main}}{}
\graphicspath{{\subfix{../../figures/}}}

\begin{document}

    \ifSubfilesClassLoaded{\tableofcontents}{}

\chapter{Results and Discussion}

Following the experimental and analysis procedures explained in\CCgen{capitolo 3}{manca sia il bersaglio sia la sintesi}{Due assenze speculari. In apertura non e' detto \emph{a che cosa servono} questi numeri: ``target'' compare una volta in tutto il capitolo, ``goal'' mai, ``requirement'' due volte, quindi il lettore non sa mai se un risultato sia buono, sufficiente o deludente. In chiusura non c'e' nessuna sezione di sintesi o discussione: il capitolo finisce sull'ultima sottosezione delle simulazioni. Il titolo dice ``Results and \emph{Discussion}'', ma la discussione vera e propria e' rimandata alle conclusioni, che pero' sono un altro capitolo e hanno un altro compito. Aggiungerei in apertura una riga per risultato che dica il bersaglio, e in chiusura mezza pagina che tiri le somme rispetto a quei bersagli.}
Chapter~\ref{cap2}, the temperature of the atomic cloud and the ODT
capture efficiency were characterized and optimized, and the stability
of the optical field delivered through the black fiber was assessed.
This chapter is organized in four parts.

The first part addresses the optimization of the molasses temperature,
carried out by adjusting the cooling-laser detuning, the cooling-laser
intensity, and the compensation magnetic field (Sec.~\ref{sec:molasses_temperature_optimization}).

The second part describes the optimization of the ODT loading. The
compensation coils are used to displace the MOT after loading, in
order to maximize the spatial overlap between the atomic cloud and the
ODT beam and thereby enhance the capture efficiency; the measured
loading efficiency is then compared with the capture model of
Chapter~\ref{cap2} (Sec.~\ref{sec:odt_loading_optimization}).

The third part presents the characterization of the power and
polarization stability of the field transmitted through the black
fiber, analyzed in the time domain with the Allan deviation and in the
frequency domain with the amplitude spectral density, and discusses the
implications of the measured noise for the optical conveyor belt
(Sec.~\ref{sec:power_polarization_results}).

The fourth part reports the numerical simulation of EIT cooling,
starting from the ideal three-level system and extending to the full
24-level structure of the $^{87}\mathrm{Rb}$ $D_2$ line, including the
effect of optical pumping and repumping on the cooling dynamics
(Sec.~\ref{sec:eit_cooling_results}).

    \import{./}{cap3_MOL_TEMP.tex}
    \clearpage
    \import{./}{cap3_ODT_LOAD.tex}
    \clearpage
    \import{./}{cap3_POLARIZATION.tex}
    \clearpage
    \import{./}{cap3_EIT.tex}


    \ifSubfilesClassLoaded{\printbibliography}{}
\end{document}